% Dirichlet Round 03 English translation TeX
% Source: G. Lejeune Dirichlet's Werke, Band II (1897), title p. 21 and body pp. 23--26.
% German source TeX prepared source-checked from dirichlet_werke_vol2_1897.pdf, PDF pages 34, 36--39.
% Prepared 2026-06-02. Faithful translation, with mathematical notation preserved.
\documentclass[11pt]{article}
\usepackage[a4paper,margin=1in]{geometry}
\usepackage[T1]{fontenc}
\usepackage[utf8]{inputenc}
\usepackage{lmodern}
\usepackage{amsmath,amssymb}
\usepackage{microtype}
\usepackage{hyperref}
\hypersetup{colorlinks=true,linkcolor=black,urlcolor=black,citecolor=black}
\setlength{\parindent}{1.2em}
\setlength{\parskip}{0.2em}
\begin{document}

\begin{center}
{\Large G. Lejeune Dirichlet's Werke. II.}\par\vspace{1.5em}
{\Large IV.}\par\vspace{0.5em}
{\LARGE \textbf{On the Reduction}}\par
{\LARGE \textbf{of Positive Quadratic Forms}}\par
{\LARGE \textbf{with Three Indeterminate Integers.}}\par\vspace{1em}
{\large By G. Lejeune Dirichlet.}\par\vspace{0.5em}
{\small Report on the proceedings of the Royal Prussian Academy of Sciences, year 1848, pp. 285--288.}
\end{center}

\begin{center}
\textbf{On}\par
\textbf{the reduction of positive quadratic forms}\par
\textbf{with three indeterminate integers.}
\end{center}

It is well known that Lagrange was the first to show that every binary quadratic form can be reduced, that is, transformed into another one whose coefficients satisfy certain inequality conditions; he also showed at the same time that in every class of positive forms there always exists only a single such form, so that in this case the various reduced forms corresponding to a given determinant can serve as the representatives of the various classes. After the ternary forms had later been considered in the \emph{Disquisitiones arithmeticae} from a general point of view, it became necessary for the further development of this theory to extend Lagrange's investigation to the positive forms of this kind, that is, to find such inequality conditions among the coefficients that, in every class, they are satisfied by one and only one form. This extension, which involved great difficulties, was accomplished by Seeber in a work devoted especially to the positive ternary forms, of which it forms the principal content. The great complication of Seeber's method made it desirable to find a simpler path leading to the same results; this path, however, could be found only after many fruitless attempts.

For an easier description of this new procedure, which is of surprising simplicity, it will be expedient to clothe the matter in a geometric form,

using here the interesting geometric construction by which Gauss, in a short essay in which he discusses Seeber's work,\footnote{\emph{Crelle's Journal}, vol. 20, p. 312. Editorial note in the \emph{Werke}: Gauss' Werke, vol. II, p. 188. K.} represents the principal properties of positive binary and ternary forms. According to this geometric representation, every positive ternary form corresponds to an infinite system of points arranged parallelepipedally, and the various subformations of the form are nothing other than changed arrangements of the same system according to another elementary parallelepiped. Expressed in this language, Seeber's results consist essentially in the following.

\begin{enumerate}
\item Every system of parallelepipedally arranged points can always be divided in such a way that, in the corresponding elementary parallelepiped, neither the sides of the faces are greater than their diagonals, nor the edges of the parallelepiped greater than the diagonals of the parallelepiped.
\item For a given system, such an arrangement can in general be carried out in only one way.
\end{enumerate}

One is easily convinced of the correctness of the first of these propositions by the following very simple consideration.

Let $(0)$ be any point of the system. The remaining points of the system evidently always occur in pairs at equal distance and in opposite directions from $(0)$. Let $(1)$ be one of the points in the pair for which the distance from $(0)$ is smaller than for every other pair. If the same least distance occurs for several pairs, choose $(1)$ arbitrarily in any one of them. Now place through the line $(01)$ and any point of the system outside that line an infinite plane; all points falling in this plane then form a parallelogrammatic system. In one of the two nearest parallel lines of this system take the point lying nearest to $(0)$, or arbitrarily one of them if the same shortest distance occurs for two points of this line. Among all planes that can be placed through $(01)$ in the indicated way, this shortest distance will be smaller in one plane, or in some of them, than in all the others. Fix this plane, or one of them if the same shortest distance should be common to several, and the point in it which we shall denote by $(2)$. Then one evidently has:
\[
(02) \geq (01),
\]
and just as easily one sees that, in the parallelogram determined by the sides $(01)$ and $(02)$, the two diagonals are not smaller than these sides. Once the plane $(012)$ has been fixed in this way, observe that the entire system of points lies in the plane $(012)$ and in other planes parallel to it and mutually equidistant. Now take, in one of the two planes adjacent to the plane $(012)$, the point nearest to $(0)$, or one of the nearest points if the same shortest distance should occur for several. If this point is called $(3)$, then one evidently has:
\[
(03) \geq (02),
\]
and the fundamental parallelepiped determined by the edges $(01)$, $(02)$, and $(03)$ will have the required properties.

After what was remarked above about the parallelogram with sides $(01)$ and $(02)$, it remains for us to show that, both in the parallelograms $(013)$ and $(023)$ and in the parallelepiped, the sides and edges are not greater than the diagonals. In view of the two inequalities
\[
(03) \geq (02) \geq (01)
\]
this evidently amounts to proving that $(03)$ is no greater than any one of the four diagonals of the named parallelograms, nor greater than any one of the four diagonals of the parallelepiped. But these eight diagonals evidently coincide, in length, with the eight straight lines that can be drawn from $(0)$ to the eight corners of the parallelograms that lie around $(3)$ in the plane through $(3)$ parallel to $(012)$. That none of these eight joining lines is smaller than $(03)$ follows immediately from the condition by which the point $(3)$ was chosen.

If, in the construction just indicated, $(1)$, $(2)$, and $(3)$ are completely determined points (apart from the possibility, always present, that for each of them one can also take the point that lies, with respect to $(0)$, at equal distance and in the opposite direction), then the edges $(01)$, $(02)$, and $(03)$ of the fundamental parallelepiped obtained are unequal, and the diagonals are actually greater than the sides and edges by which, in general, they are not to be exceeded. In this case one then easily proves

that this system admits only this single elementary parallelepiped with the required properties.

The matter takes a somewhat different form when, in our construction, a choice can be made among different non-opposite positions for the points $(1)$, $(2)$, and $(3)$. In such singular cases there may exist, for the system, several non-congruent fundamental parallelepipeds possessing the required properties. All these parallelepipeds, however, can be completely enumerated without difficulty, and one can always distinguish one of them from the others by certain subsidiary conditions, so that, with the addition of these secondary conditions, the theorem that every class contains only one reduced form does not lose its validity. Something similar is already well known in the theory of positive binary forms, where in two special cases one must prescribe a definite sign for the middle coefficient if one wishes to retain only one reduced form in each class.

\end{document}
